\input{preamble.tex}

\title{Лабораторная работа № 4 \\ Первоначальное конфигурирование сети}
\author{Сущенко Алина \\ НПИбд-01-23}

\institute{Российский университет дружбы народов имени Патриса Лумумбы}
\date{2026}

\begin{document}

\frame{\titlepage}

\begin{frame}
\frametitle{Цель работы}
\begin{itemize}
    \item Провести подготовительную работу по первоначальной настройке коммутаторов сети
\end{itemize}
\end{frame}

\begin{frame}
\frametitle{Схема подключения оборудования}
    \centering
    \includegraphics[width=0.9\textwidth]{../images/01.png}
\end{frame}

\begin{frame}
    \centering
    \includegraphics[width=0.8\textwidth]{../images/02.png}
\end{frame}

\begin{frame}
    \centering
    \includegraphics[width=0.8\textwidth]{../images/03.png}
\end{frame}

\begin{frame}
    \centering
    \includegraphics[width=0.8\textwidth]{../images/04.png}
\end{frame}

\begin{frame}
    \centering
    \includegraphics[width=0.8\textwidth]{../images/05.png}
\end{frame}

\begin{frame}
    \centering
    \includegraphics[width=0.8\textwidth]{../images/06.png}
\end{frame}

\begin{frame}
    \centering
    \includegraphics[width=0.8\textwidth]{../images/07.png}
\end{frame}

\begin{frame}
    \centering
    \includegraphics[width=0.8\textwidth]{../images/08.png}
\end{frame}

\begin{frame}
    \centering
    \includegraphics[width=0.8\textwidth]{../images/09.png}
\end{frame}

\begin{frame}
    \centering
    \includegraphics[width=0.8\textwidth]{../images/10.png}
\end{frame}

\begin{frame}
    \centering
    \includegraphics[width=0.8\textwidth]{../images/11.png}
\end{frame}

\begin{frame}
    \centering
    \includegraphics[width=0.8\textwidth]{../images/12.png}
\end{frame}

\begin{frame}
    \centering
    \includegraphics[width=0.8\textwidth]{../images/13.png}
\end{frame}

\begin{frame}
    \centering
    \includegraphics[width=0.8\textwidth]{../images/14.png}
\end{frame}

\begin{frame}
    \centering
    \includegraphics[width=0.8\textwidth]{../images/15.png}
\end{frame}

\begin{frame}
    \centering
    \includegraphics[width=0.8\textwidth]{../images/16.png}
\end{frame}

\begin{frame}
    \centering
    \includegraphics[width=0.8\textwidth]{../images/17.png}
\end{frame}

\begin{frame}
\frametitle{Сводная информация по настройкам}
В процессе выполнения работы все коммутаторы были настроены в соответствии с планом IP-адресации:

\begin{table}[h]
    \centering
    \begin{tabular}{|c|c|c|}
        \hline
        \textbf{Устройство} & \textbf{Имя} & \textbf{IP-адрес} \\
        \hline
        Коммутатор 1 & msk-ansusenko-sw-01 & 10.128.1.2 \\
        \hline
        Коммутатор 2 & msk-ansusenko-sw-1 & 10.128.1.6 \\
        \hline
        Коммутатор 3 & msk-ansusenko-sw-2 & 10.128.1.3 \\
        \hline
        Коммутатор 4 & msk-ansusenko-sw-3 & 10.128.1.4 \\
        \hline
        Коммутатор 4 & msk-ansusenko-sw-4 & 10.128.1.5 \\
        \hline
    \end{tabular}
\end{table}
\end{frame}

\begin{frame}
\frametitle{Контрольные вопросы}

\begin{block}{Вопрос 1: При помощи каких команд можно посмотреть конфигурацию сетевого оборудования?}
Для просмотра конфигурации сетевого оборудования Cisco используются следующие команды:
\begin{itemize}
    \item \texttt{show running-config} (или \texttt{show run}) — отображает текущую конфигурацию
    \item \texttt{show startup-config} (или \texttt{show start}) — отображает стартовую конфигурацию
    \item \texttt{show ip interface brief} — показывает информацию по интерфейсам
    \item \texttt{show vlan} — отображает информацию о VLAN
    \item \texttt{show interfaces trunk} — показывает информацию о транковых портах
\end{itemize}
\end{block}
\end{frame}

\begin{frame}
\frametitle{Контрольные вопросы}

\begin{block}{Вопрос 2: При помощи каких команд можно посмотреть стартовый конфигурационный файл оборудования?}
Стартовый конфигурационный файл (startup-config) просматривается командой:
\begin{itemize}
    \item \texttt{show startup-config}
    \item \texttt{show start} (сокращенный вариант)
\end{itemize}
\end{block}
\end{frame}

\begin{frame}
\frametitle{Контрольные вопросы}

\begin{block}{Вопрос 3: При помощи каких команд можно экспортировать конфигурационный файл оборудования?}
Для экспорта конфигурации используются следующие команды:
\begin{itemize}
    \item \texttt{write memory}
    \item \texttt{copy running-config startup-config}
    \item \texttt{copy running-config tftp:}
    \item \texttt{copy startup-config tftp:}
\end{itemize}
\end{block}
\end{frame}

\begin{frame}
\frametitle{Контрольные вопросы}

\begin{block}{Вопрос 4: При помощи каких команд можно импортировать конфигурационный файл оборудования?}
Для импорта конфигурации используются следующие команды:
\begin{itemize}
    \item \texttt{copy startup-config running-config}
    \item \texttt{copy tftp: running-config}
    \item \texttt{copy tftp: startup-config}
    \item \texttt{reload}
\end{itemize}
\end{block}
\end{frame}

\begin{frame}
\frametitle{Выводы}
\begin{itemize}
    \item В результате выполнения работы мы успешно воссоздали схему подключения оборудования, настроили её и провели проверку текущих настроек.
\end{itemize}
\end{frame}

\end{document}
